\documentclass[t,12pt]{beamer}

\usepackage{amssymb}
\usepackage{amsfonts}
\usepackage{amsmath}
\usepackage{dsfont}
\usepackage[utf8x]{inputenc}
\usepackage[labelfont=bf,justification=justified]{caption}
\usepackage{subcaption}
\usepackage{fixltx2e}
\usepackage{setspace}
\usepackage{arydshln}
\usepackage{xcolor}
\usepackage{tikz}
\usepackage{chngpage}
\usepackage{comment}
\usetikzlibrary{shapes,arrows}
\tikzstyle{block} = [rectangle, draw, fill=blue!20, 
    text width=5em, text centered, rounded corners, minimum height=4em]
\tikzstyle{line} = [draw, -latex']

\usetheme{Emory}
\makeatletter
\setbeamertemplate{footline}
{
  \leavevmode%
  \hbox{%
  \begin{beamercolorbox}[wd=.333333\paperwidth,ht=2.25ex,dp=1ex,center]{author in head/foot}%
    \usebeamerfont{author in head/foot}\insertshortauthor~~\beamer@ifempty{\insertshortinstitute}{}{\insertshortinstitute}
  \end{beamercolorbox}%
  \begin{beamercolorbox}[wd=.333333\paperwidth,ht=2.25ex,dp=1ex,center]{title in head/foot}%
    \usebeamerfont{title in head/foot}\insertshorttitle
  \end{beamercolorbox}%
  \begin{beamercolorbox}[wd=.333333\paperwidth,ht=2.25ex,dp=1ex,right]{date in head/foot}%
    \usebeamerfont{date in head/foot}\insertshortdate{}\hspace*{2em}
    \insertframenumber{} / \inserttotalframenumber\hspace*{2ex} 
  \end{beamercolorbox}}%
  \vskip0pt%
}
\makeatother

\setbeamerfont{frametitle}{size=\large}
\setbeamertemplate{navigation symbols}{}

\titlegraphic{\includegraphics[width=0.3\textwidth]{~/Dropbox/Emory.png}}

\begin{document}

\title[Exponential smoothing]{Exponential smoothing}
\author{\vspace{0.35in} \\ {\fontfamily{ppl}\selectfont William Mann} \vspace{-0.35in}}

\date{}
\maketitle

	\begin{frame}[c]
	\vfill
	\centerline{Simple exponential smoothing}
	\vfill
	\end{frame}

	\begin{frame}[c]{Simple exponential smoothing}
	Exponential smoothing is a family of weighted-average methods.
	\\
	The simplest example is called \textbf{simple} exponential smoothing:
	\begin{itemize}
	\medskip
	\item First you make a forecast $F_1$ about the very first observation $y_1$.
	\smallskip
	\item After you observe $y_1$, your forecast of $F_2$ is given by
	$$
	F_2 = \alpha y_1 + (1-\alpha) F_1
	$$
	\item Similarly, after you observe $y_2$, your next forecast is
	$$
	F_3 = \alpha y_2 + (1-\alpha) F_2
	$$
	and so forth.
	\smallskip
	\item The value of $\alpha$ is called a smoothing parameter.
	\end{itemize}
	\end{frame}

	\begin{frame}[c]{Why is it called exponential smoothing?}
	The procedure on the previous slide gives us this formula:
	$$
	F_{t+1} = \alpha \times \left[ y_t + (1-\alpha)y_{t-1} + (1-\alpha)^2y_{t-2} + ... \right]
	$$
	As time goes by, your initial forecast $F_1$ no longer matters.
	\\ \bigskip
	There is always much greater weight given to recent observations rather than past observations, but at the same time,
	the process ``remembers'' a very long history of information about $y_t$.
	\\ \bigskip
	This approach achieves a lot for the cost of only one parameter $\alpha$.
	\end{frame}

	\begin{frame}[c]{Implementation: Choice of initial forecast}
	How to choose the initial forecast $F_1$? This should not be a big issue.
	\begin{itemize}
	\medskip
	\smallskip
	\item Any different choice of $F_1$ will give you slightly different results, but it should not be a big issue in practice.
	\smallskip
	\item If two reasonable choices yield very different results, something may be going wrong, or you may just need more data.
	\item The simplest choice is to ``look ahead'' and just use $y_1$.
	\\ This is what we will do when implementing by hand.
	\smallskip
	\item In this case the first forecast is slightly cheating, and we should ignore it when assessing the performance of the model.
	\smallskip
	\item Software will often ``backcast'' from the early observations.
	\\ Every common approach involves looking ahead in some way.
	\end{itemize}
	\end{frame}

	\begin{frame}[c]{Implementation: Choice of smoothing parameter $\alpha$}
	How to choose $\alpha$? This is a more important issue.
	\begin{itemize}
	\medskip
	\item Notice that $\alpha$ is the amount that we change our forecast in response to a new observation.
	So actually a higher $\alpha$ makes our forecasts \emph{less} smooth.
	(This is a bit confusing but is standard!)
	\smallskip
	\item If the goal is just to visualize data, it's common to just try values of $\alpha$ until one looks reasonable.
	\smallskip
	\item If the goal is to forecast, then $\alpha$ is typically chosen to minimize the sum of squared one-step forecast errors,
	$
	\sum_t (y_t - F_t)^2
	$.
	\smallskip
	\item 
	Although it looks much like a regression, there is no simple formula for this.
	Instead, search numerically over values of $\alpha$.
	\end{itemize}
	\medskip
	See examples on Canvas.
	\end{frame}

	\begin{frame}[c]
	\vfill
	\centerline{Double exponential smoothing}
	\centerline{(also called Holt's method)}
	\vfill
	\end{frame}

	\begin{frame}[c]{Double exponential smoothing}
	When a series is trending in one direction, simple exponential smoothing will always generate forecasts that lag behind.
	\\ \medskip
	Double exponential smoothing keeps track of both the \emph{level} of the series (as before), and also its \emph{trend}.
	\\ \medskip
	The second part gives a sense of where the series is going.
	This improves forecast accuracy for data with slowly-evolving trends.
	\\ \medskip
	The formal framework is as follows:
	\begin{align*}
	L_t \quad &= \quad \alpha Y_t \, + \, (1-\alpha)(L_{t-1} + T_{t-1})
	\\
	T_t \quad &= \quad \beta (L_t - L_{t-1}) \, + \, (1-\beta) T_{t-1}
	\\
	F_{t+1} \quad &= \quad L_t \, + \, T_t
	\end{align*}
	{\footnotesize (There is also a popular version with multiplicative trends, but it is easier and almost identical to just log-transform the data in the above approach.)}
	\end{frame}

	\begin{frame}[c]{Double exponential smoothing: Implementation}
	How to set the initial values of the forecast components, $L_1$ and $T_1$?
	\begin{itemize}
	\item As before, we can set $L_1 = Y_1$. 
	For $T_1$, a common choice is $T_1 = Y_2 - Y_1$, and this is what we will do in our examples. 
	\item Most software follows an approach based on fitting a simple regression model to the early data.
	\item As with SES, this choice generally should not make a very big difference to your forecasts.
	\end{itemize}
	\medskip
	How to choose the parameters $\alpha$ and $\beta$?
	\begin{itemize}
	\item Again, just minimize the sum of squared one-step forecast errors.
	\item We just have to search over two parameters instead of one.
	\end{itemize}
	\end{frame}

	\begin{frame}[c]
	\vfill
	\centerline{Triple exponential smoothing}
	\centerline{(also called the Holt-Winters approach)}
	\vfill
	\end{frame}

	\begin{frame}[c]{Triple exponential smoothing, aka Holt-Winters}
	Simple exponential smoothing captures the current \emph{level} of a series. 
	\\
	Double exponential smoothing learns which way it's \emph{trending}.
	\\ 
	The other big component of most economic data is \emph{seasonality}.
	\\ \bigskip
	Neither simple nor double smoothing can capture seasonality,
	\\ but we can extend the method one more time to do so.
	\\ \bigskip
	The result is called \textbf{triple exponential smoothing} or the \textbf{Holt-Winters method}, or sometimes just \textbf{Winters method}.
	\\ \bigskip
	This has been an extremely popular and effective forecasting method for many decades now.
	\end{frame}

	\begin{frame}[c]{Triple exponential smoothing: Framework}
	\vspace{-0.4in}
	\begin{align*}
	L_t \quad &= \quad \alpha (Y_t-S_{t-p}) \,\, + \,\, (1-\alpha)(L_{t-1} + T_{t-1})
	\\
	T_t \quad &= \quad \beta (L_t - L_{t-1}) \,\, + \,\, (1-\beta) T_{t-1}
	\\
	S_t \quad &= \quad \delta (Y_t - L_t) \,\, + \,\, (1-\delta) S_{t-p}
	\\
	F_t \quad &= \quad L_{t-1} \,\, + \,\, T_{t-1} \,\, + \,\, S_{t-p}
	\end{align*}
	We will not try to implement this by hand. 
	The extra math, and the question of how exactly to initialize $S_t$, are tricky but not interesting.
	\\ \medskip
	The main takeaway is that this approach can capture rich behavior, 
	\\ and can quickly generate accurate forecasts in most practical settings,
	\\ while estimating just three parameters ($\alpha$, $\beta$, $\delta$).
	\\ \smallskip
	{\footnotesize
	Again, there is also a version with \emph{multiplicative} seasonality, 
	in which we divide and multiply by $S_t$ and $S_{t-p}$, instead of subtracting and adding.
	For some reason this version is very popular.
	In my experience the two versions give similar results.}
	\end{frame}

	\begin{frame}[c]{Exponential smoothing: Summary}
	\begin{itemize}
	\item
	Exponential-smoothing is a very popular forecasting approach.
	\\ But notice that we derived it in a very ad-hoc way.
	\smallskip
	\item
	As long as we are \emph{just} forecasting, this may not be a problem.
	\smallskip
	\item
	But there is less clear guidance for other goals, such as building prediction intervals, connecting with economic models, etc.
	\smallskip
	\item
	These questions eventually require us to model the \emph{process} that generates the data, as is usually done in statistics.
	\smallskip
	\item
	From this point forward, we will develop such a framework.
	\\ It is often called the ARMA framework, for reasons we will see.
	\smallskip
	\item
	The point is not to generate better forecasts (it won't), 
	\\ but rather to give a clearer understanding of what we are doing.
	\end{itemize}
	\end{frame}

\end{document}
